%% sfc_appendix.tex — Supplementary material (Appendices A–D) for sfc.tex.
%% Standalone document: compiles independently of the main paper.
%%   pdflatex sfc_appendix
%%   bibtex   sfc_appendix
%%   pdflatex sfc_appendix
%%   pdflatex sfc_appendix

\documentclass[journal]{IEEEtran}
\IEEEoverridecommandlockouts

\usepackage{cite}
\usepackage{amsmath,amssymb,amsfonts}
\usepackage{graphicx}
\usepackage{booktabs}
\usepackage{xcolor}
\usepackage{url}
\usepackage[colorlinks=true,linkcolor=blue,citecolor=blue,urlcolor=blue]{hyperref}

%% ── Notation macros (kept in sync with sfc.tex) ─────────────────────────────
\newcommand{\PORTER}{\textsc{PORTER}}
\newcommand{\TeleSABRE}{\textsc{TeleSABRE}}
\newcommand{\dSABRE}{\textsc{dSABRE}}
\newcommand{\SABRE}{\textsc{SABRE}}
\newcommand{\TeleData}{\textsc{TeleData}}
\newcommand{\TeleGate}{\textsc{TeleGate}}
\newcommand{\TeleGatePacket}{\textsc{TeleGatePacket}}
\newcommand{\logi}[1]{\mathrm{q}_{#1}}
\newcommand{\phys}[1]{\mathrm{p}_{#1}}

%% ── References into the main paper ──────────────────────────────────────────
%% The supplement compiles independently, so these are hardcoded; update here
%% (single place) if the main paper is renumbered.
\newcommand{\mainInitialLayout}{Section~III-A}   % sec:sfc-layout
\newcommand{\mainRunningExample}{Section~III-G}  % sec:sfc-example
\newcommand{\mainExampleFig}{Figure~5}           % fig:example
\newcommand{\mainEvaluation}{Section~IV}         % sec:experiments
\newcommand{\mainSetup}{Section~IV-B}            % sec:setup
\newcommand{\mainProtocols}{Section~IV-C}        % sec:protocols
\newcommand{\mainSharedResults}{Section~IV-E}    % sec:shared_results
\newcommand{\mainCompileTime}{Section~IV-G}      % sec:time
\newcommand{\mainSharedTable}{Table~VII}         % tab:shared

\begin{document}

\title{Supplementary Material for\\
  ``PORTER: A Plan-Committed Router for Multi-Core Distributed Quantum
  Computers''}

%\author{Fuchao~Sheng, xxx, xxx, Sanjiang~Li}

\maketitle

\noindent
This supplement contains Appendices A--D of the main paper. Section,
figure, and table numbers of the form ``Section~III-A'' or
``Table~VII'' refer to the main paper; numbers local to this document
refer to the supplement itself.

\appendices

\section{The 80-Qubit Suite: Reconstruction Detail}
\label{app:suite80}

The 80-qubit suite (\mainSetup{} of the main paper) was generated with
\texttt{mqt.bench}~2.2.2~\cite{quetschlich2023mqtbench}
(device-independent generation followed by
transpilation to our basis, since the installed \texttt{mqt.bench}'s
target-aware native-gate path raised an unrelated \texttt{qiskit}
compatibility error on every circuit we tried). AE at this size is not
directly producible by stock \texttt{mqt.bench}: its phase-estimation
construction evaluates $\mathrm{Operator}(Q)^{2^{j}}$ for each
evaluation bit $j$, and the accumulated floating-point non-unitarity
trips \texttt{qiskit}'s unitarity check
(\texttt{"Input matrix is not unitary"}) beyond roughly $27$ evaluation
qubits, which is why published AE files stop near $54$~qubits. Since the
Bernoulli operator $Q$ is a single-qubit rotation, we instead supply the
power analytically, $Q^{2^{j}}=R_Y(2^{j+1}\theta_p)$ ($R_Y$ periodic in
$4\pi$), and otherwise run \texttt{mqt.bench}'s own native-gate
transpilation unchanged. The CNOT sub-sequence---which alone determines
routing---is decomposition-invariant in these rotation angles and is
byte-identical to stock \texttt{mqt.bench} at every size the latter can
build (verified at $10$, $20$, $25$, and $28$~qubits); only the
routing-irrelevant single-qubit angles carry floating-point error at
high bit indices.

QNN required an analogous reconstruction. Between
\texttt{mqt.bench}~1.1.0 (which produced the deep 8{,}126-CX 64q QNN
this suite uses) and the installed 2.2.2, the benchmark's definition
was changed from a ZZ feature map (full all-pairs entanglement) plus a
\texttt{RealAmplitudes} ansatz to a plain Z feature map (no
entanglement) plus the same ansatz. The current generator therefore
emits a near-linear 79-CX circuit at 80 qubits whose two-qubit
interaction graph is \emph{identical} to GHZ's --- a degenerate
duplicate, not a distinct benchmark. We restore the original deep QNN
by rebuilding the 1.1.0 construction (ZZ feature map, \texttt{reps}=2,
full entanglement; \texttt{RealAmplitudes}, \texttt{reps}=1) and
transpiling it to our basis; at 64 qubits this reproduces the checked-in
1.1.0 file's CNOT sub-sequence byte-for-byte (identical ordered pair
sequence, all $\binom{64}{2}$ interacting pairs, 8{,}126 CX), so the
80-qubit instance (a fully-connected $\binom{80}{2}$-edge interaction
graph, 12{,}718 CX) is a faithful extrapolation of the same
construction.

\section{External Layouts from pytket-dqc}
\label{app:pytket}

As a third, independent layout source complementing the shared-layout
protocol of \mainProtocols{} of the main paper, we route \dSABRE{} and
\PORTER{} from a placement generated by \texttt{pytket-dqc}'s
\texttt{HypergraphPartitioning} allocator~\cite{pytketdqc2024}
(best-of-3 KaHyPar seeds on a \texttt{NISQNetwork} matching our core
graph), converted to a physical layout by the same topology-aware
intra-core assignment \dSABRE{}'s own pytket benchmark
uses~\cite{lidsabre}: each core's assigned logical qubits are placed
by descending interaction-graph degree onto physical positions ordered
by ascending intra-core centrality, then locally refined by a bounded
pairwise-swap pass. \TeleSABRE{} is excluded for the same structural
reason as in \mainProtocols{} (no externally supplied
layout support in its CLI). \texttt{HypergraphPartitioning} only
decides which \emph{core} each logical qubit goes to; the capacity it
sees for a core is the length of that core's declared qubit list in
the \texttt{NISQNetwork}. Following \dSABRE{}'s stock corner-removal
convention (the four lowest-degree corner qubits of the whole-chip
graph; \mainProtocols{}), we withhold those four corner
qubits from the cores' declared capacity and from the later
physical-placement step within each core. This guarantees only a
four-qubit margin \emph{chip-wide}, not per-core, so a core can still
end up fully packed; rather than skip the run in that case,
\PORTER{}'s general safeguard against fully packed cores
(\mainInitialLayout{}) legalizes the layout automatically
before routing begins. \dSABRE{}'s teledata-only mechanism does not
need a free port on the side it teleports \emph{from} and was never
affected by this precondition. Because \texttt{pytket-dqc} uses gate
teleportation, not state teleportation, its own reported cost is in a
different unit (e-bits per Steiner-tree-amortised hyperedge, not EPR
pairs per qubit move) and is included only as context, following
\dSABRE{}'s own established convention for this
comparison~\cite{lidsabre}.

\begin{table}[t]
\centering
\caption{Pytket-mapping shared-layout comparison (25q and 64q suites),
  corners-removed \texttt{NISQNetwork}
  capacity. \texttt{pytket e-bits} is
  \texttt{HypergraphPartitioning}'s own cost in a different unit
  (context only, not EPR-comparable). \PORTER{}'s general full-core
  safeguard leaves no missing entries on
  this suite.}
\label{tab:pytket}
\setlength{\tabcolsep}{3pt}
\small
\resizebox{\columnwidth}{!}{%
\begin{tabular}{@{}l l r rr rr@{}}
\toprule
Suite & Circuit & \texttt{pytket} e-bits & \multicolumn{2}{c}{\dSABRE{}} & \multicolumn{2}{c}{\PORTER{}} \\
\cmidrule(lr){4-5}\cmidrule(lr){6-7}
& & & EPR & SWAP & EPR & SWAP \\
\midrule
25q & ae         &  78 &   67 &  339 &   73 &  348 \\
25q & ghz        &   3 &    3 &   14 &    3 &   17 \\
25q & graphstate &   4 &    5 &   24 &    4 &   19 \\
25q & qft        &  70 &   92 &  437 &  118 &  438 \\
25q & qnn        & 148 &  168 &  848 &  127 &  703 \\
25q & random     & 564 &  440 & 1482 &  395 & 1340 \\
64q & ae         & 264 &  361 & 2243 &  223 & 1638 \\
64q & ghz        &   5 &    5 &   74 &    5 &   62 \\
64q & graphstate &  10 &    8 &   68 &   10 &   60 \\
64q & qft        & 307 &  320 & 2112 &  278 & 2021 \\
64q & qnn        & 621 &  509 & 8215 & 1031 & 6538 \\
64q & random     & 1020 & 750 & 3969 &  861 & 4384 \\
\bottomrule
\end{tabular}%
}
\end{table}

Table~\ref{tab:pytket} reports
a complete grid: \PORTER{}'s full-core safeguard
legalizes the three layouts that would
otherwise violate its teleport precondition (25q QNN, 64q AE,
64q QNN), at the cost of the layout-quality hit an emergency eviction
inevitably takes -- most visibly on 64q QNN, where the evicted layout
lands \PORTER{} at 1{,}031 EPR pairs against \dSABRE{}'s 509, its worst
relative result on this suite. Excluding that outlier, results split
close to evenly: \PORTER{} uses fewer EPR pairs on 25q
Graphstate/QNN/Random and 64q AE/QFT, \dSABRE{} on 25q AE/QFT and 64q
Graphstate/Random, and the two tie on both GHZ rows; geometric-mean
EPR across all twelve rows is within $1.4\%$ of parity
(\dSABRE{} $69.2$ vs.\ \PORTER{} $70.2$). Pooled with the
\TeleSABRE{}-layout grid of \mainSharedTable{} of the main paper,
\PORTER{} uses
fewer EPR pairs on fifteen of the thirty fixed-layout instances,
\dSABRE{} on twelve, with three ties
(\mainSharedResults{}).

The grid above feeds pytket-dqc's partition to \PORTER{} as a
controlled probe. The complementary pipeline question -- offering the
same partition to \PORTER{} as an optional extra layout candidate
beside its own, which lowers geometric-mean EPR by a further $16\%$ on
the most partitionable circuits but is kept opt-in on cost and
robustness grounds -- is a configuration study of \PORTER{} itself
rather than part of this comparison, and is deferred to
Appendix~\ref{app:optin}.

\section{pytket-dqc as an Opt-In \PORTER{} Layout Candidate}
\label{app:optin}

Appendix~\ref{app:pytket} feeds pytket-dqc's partition to
\PORTER{} as a controlled probe; the same partition can also be
offered to \PORTER{} as one extra layout candidate beside its own
(\mainInitialLayout{} of the main paper). Since
\PORTER{} routes far faster than pytket-dqc partitions, this is nearly
free on the routing side; taking the better of the two per instance
lowers the eighteen-instance geometric-mean EPR from $116.8$ to $98.6$
($-16\%$), with the gains concentrated on the most partitionable
circuits (Graphstate $7\!\to\!4$, $20\!\to\!10$, $28\!\to\!11$ across
the three suites; 80q~AE $490\!\to\!352$). We keep it \emph{opt-in}
rather than the default for three reasons. First, on the dense and
Random circuits the external partition is often \emph{worse} for
\PORTER{} (64q~QNN $864\!\to\!1{,}031$, 64q~Random $737\!\to\!861$), so
the best-of simply retains \PORTER{}'s own layout there. Second, the
partition itself is expensive -- seconds to minutes, and on dense
circuits more than \PORTER{}'s entire routing pass ($309$\,s to
partition 80q~Random versus $\approx\!37$\,s to route it; $16$\,s
versus $0.6$\,s on 80q~AE) -- against a sub-millisecond cost for
\PORTER{}'s own layout build, so adopting it by default would erase
\PORTER{}'s compile-time advantage (\mainCompileTime{} of the main
paper). Third, it
reintroduces the external-partitioner dependency \PORTER{}'s
self-contained layout avoids, and unlike it can leave a core fully
packed, relying on the full-core repair of
\mainInitialLayout{}. \PORTER{} therefore defaults to its own
layout and treats an external partition as an optional refinement,
selected only when EPR is at a premium and its generation time is
acceptable.

\section{Independent Replay Verification}
\label{app:verification}

Beyond code review, we built an independent replay verifier that
treats the router as a black box: it takes only the emitted
\texttt{CompiledOp} stream, the original circuit, and the initial
layout, and re-derives, from scratch, whether the stream is a faithful
and physically legal execution. The verifier maintains its own
physical-position map and a per-logical-qubit FIFO queue of pending
gate IDs (derived directly from the circuit, independent of any
router-internal DAG state), and checks, per emitted operation:
\begin{itemize}
\item every gate is popped from the head of \emph{every} one of its
  operand qubits' queues before being marked executed, so gates cannot
  execute out of DAG order or twice;
\item every \texttt{local\_swap}/\texttt{local\_cx} acts on a genuine
  intra-core coupling edge, at the physical positions the verifier's
  own simulated layout says the operand logical qubits occupy;
\item every \texttt{teleport} uses the correct inter-core port pair for
  its two cores, both ports were free beforehand, and the source
  qubit was coupling-adjacent to the source port;
\item every \texttt{cat\_entangler}, \texttt{cat\_cx}, and
  \texttt{cat\_disentangler} triple observes port-adjacency and
  free-port preconditions, and no other operation touches either
  communication port while a cat section is open;
\item the emitted \texttt{RouterStats} counters (EPR, teleport,
  telegate, packet, SWAP, local-CX, single counts) match the counts
  the verifier itself derives from the op stream, catching any
  bookkeeping drift between what is emitted and what is counted.
\end{itemize}
We ran the verifier against 60 randomized circuits (4--40 logical
qubits, 20--600 gates, both architectures, seeds 0--59) and against
every circuit in the evaluation suite (\mainEvaluation{} of the main
paper);
all pass. Two further artefacts make these checks concrete: a
\emph{certified trace} (Figure~\ref{fig:trace}) and a
\emph{fault-injection study} (Table~\ref{tab:faults}).

\begin{figure}[t]
\centering
\scriptsize
\setlength{\tabcolsep}{3.5pt}
\begin{tabular}{@{}r l p{4.0cm}@{}}
\toprule
\# & emitted operation & verifier asserts \\
\midrule
0 & \texttt{swap} $\phys{3}{\leftrightarrow}\phys{4}$ &
    coupling edge; operands sit where the replayed layout says \\
\midrule
1 & \texttt{cat\_entangler} $\logi{0}$; $(\phys{5},\phys{12})$; 1 EPR &
    no section already open; \emph{exact} port pair of $(c_0,c_1)$;
    both ports free; $\logi{0}$ adjacent to $\phys{5}$; opens section,
    reserving $\phys{5},\phys{12}$ \\
2 & \texttt{swap} $\phys{10}{\leftrightarrow}\phys{9}$ \hfill route $\logi{1}$ &
    tagged as packet routing; touches neither reserved port \\
3 & \texttt{cat\_cx} $g_1$ $(\phys{12},\phys{9})$ &
    $g_1$ at the head of \emph{both} $\logi{0}$'s and $\logi{1}$'s
    queues; control $=$ section control; proxy is the target-side
    port $\phys{12}$; target in $c_1$, adjacent to $\phys{12}$ \\
4 & \texttt{swap} $\phys{13}{\leftrightarrow}\phys{14}$ \hfill route $\logi{2}$ &
    as op 2 \\
5 & \texttt{cat\_cx} $g_2$ $(\phys{12},\phys{13})$ & as op 3 \\
6 & \texttt{cat\_cx} $g_3$ $(\phys{12},\phys{15})$ &
    as op 3 ($\logi{3}$ already port-adjacent) \\
7 & \texttt{cat\_disentangler} &
    a section is open; closes it, releasing the ports \\
\midrule
8 & \texttt{swap} $\phys{6}{\leftrightarrow}\phys{7}$ & coupling edge \\
9 & \texttt{local\_cx} $g_5$ $(\phys{4},\phys{7})$ &
    queue heads; same core; coupling edge \\
10 & \texttt{swap} $\phys{13}{\leftrightarrow}\phys{14}$ & coupling edge \\
11 & \texttt{teleport} $\logi{5}$: $\phys{13}{\to}\phys{5}$; 1 EPR &
    cores adjacent; exact port pair $(\phys{12},\phys{5})$; both
    ports free; $\logi{5}$ adjacent to $\phys{12}$; lands exactly on
    $\phys{5}$ \\
12 & \texttt{swap} $\phys{4}{\leftrightarrow}\phys{5}$ & coupling edge \\
13 & \texttt{local\_cx} $g_4$ $(\phys{3},\phys{4})$ &
    queue heads; coupling edge \\
\bottomrule
\end{tabular}
\caption{The full operation stream \PORTER{} emits for the running
  example of the main paper (\mainRunningExample{},
  \mainExampleFig{}), as replayed by the verifier.
  Ops 1--7 execute inside the live cat section, where only
  packet-tagged SWAPs and 1Q gates that avoid the two reserved ports
  are legal. At end of stream the verifier further asserts: no cat
  section left open, all five gates executed exactly once, and the
  emitted statistics equal to the derived counts (2 EPR, 1 teleport,
  1 packet, 0 single telegates, 6 SWAPs, 2 local CXs).}
\label{fig:trace}
\end{figure}

\begin{table}[t]
\caption{Fault-injection validation of the verifier
  (\texttt{verify\_fault\_injection.py}): single-point corruptions
  applied to operation streams that pass verification, each of which
  must be rejected.}
\label{tab:faults}
\centering
\scriptsize
\begin{tabular}{@{}ll@{}}
\toprule
Injected corruption & Rejecting check \\
\midrule
delete a \texttt{cat\_disentangler} & section state machine \\
swap a teleport's two ports & exact inter-core port pair \\
zero a teleport's EPR cost & per-op EPR accounting \\
duplicate a gate-carrying op & every gate exactly once \\
delete a gate-carrying op & queue order / end-of-stream coverage \\
swap two ops sharing a qubit & per-qubit dependency order \\
retarget a SWAP to a non-coupled pair & coupling-edge legality \\
untag an in-section routing SWAP & cat-section isolation \\
retarget an in-section SWAP onto a port & reserved-port protection \\
increment \texttt{stats.epr} by one & stats--stream cross-check \\
\bottomrule
\end{tabular}
\end{table}

\textbf{A certified trace.} Figure~\ref{fig:trace} replays the
stream \PORTER{} emits for the running example of the main paper
(\mainRunningExample{}) --- fourteen operations exercising
every op type --- and lists, per operation, exactly what the verifier
asserts before accepting it. The
trace makes the three primitives'
legality conditions operational: the \TeleGatePacket{} is one
\emph{bracketed section} in which the two reserved ports are
untouchable and every interleaved SWAP must carry packet provenance;
the \TeleData{} at op 11 must name the unique port pair of its core
pair, find both ports free, start port-adjacent, and land exactly on
the destination port; and every gate execution, local or remote, is
popped from the head of each operand's pending-gate queue, so DAG
order is enforced gate by gate rather than assumed from the router's
own bookkeeping.

\textbf{Fault injection.} On streams that always pass, a verifier
that checks nothing is indistinguishable from one that checks
everything, so we also certify that each check \emph{fires}:
\texttt{verify\_fault\_injection.py} takes passing op streams and
applies the ten single-point corruptions of Table~\ref{tab:faults},
one mutant per corruption, spanning misuse of all three communication
primitives (wrong ports, free-port violations, broken cat-section
bracketing), gate-coverage and dependency-order violations, illegal
SWAPs, and statistics drift. Across the running example (which
admits all ten corruption classes) and thirty fuzzed circuits on both
architectures, all 276 applicable mutants are rejected, each by the
intended check; the same harness runs on any benchmark QASM file.

\bibliographystyle{IEEEtran}
\bibliography{sfc}

\end{document}
